\section{Conclusion}
\label{sec:conclusion}

We asked which RIS phase optimizer a base station should run when users move, and for how long. Counting the time each method needs, we find the following.
\begin{enumerate}
    \item \textbf{Speed beats precision.} The solvers with the best rate on paper lose to a single-user heuristic that takes milliseconds, even with a tuned number of coordinate sweeps; without acquisition delay, this holds at every speed from 3 to 30~km/h (with optimistic values for the slow solvers). Averaged over longer windows, coordinate descent can tie or pass the heuristic but stays behind tuned Adam.
    \item \textbf{A few steps are optimal.} Adam's narrowband effective rate peaks after a handful of steps and then falls when the phase is renewed every few tens of milliseconds (holding it longer moves a flat optimum up); the optimum moves from 8 steps at 3~km/h to 3 at 10~km/h, the upper end of the IMT-2020 pedestrian class.
    \item \textbf{One number sets the optimum.} The best step count is governed by the step time relative to the Doppler frequency, $x=\Delta T f_d$, and scales approximately as $x^{-2/(\gamma+2)}$ (measured exponent $-0.675$); with a hundredfold faster implementation, convergence becomes nearly free.
    \item \textbf{Pilots set the update period.} In continuous operation the pilot overhead decides how often to update; each update should still be short.
    \item \textbf{Learning must win on the clock.} Winning at equal steps is not enough: the learned update rule (static set, one CPU thread) is dominated once each step is charged its measured time.
\end{enumerate}
In the regime quantified here, latency is not a secondary metric for mobile RIS: it decides which method is best.

\emph{Practical rules.} Choose the number of steps from $x$, the acquisition delay and how long the phase will be held, not from convergence. Choose the update period from the pilot overhead. Compare methods on measured latency against a classical baseline tuned at each budget, and report dominance rather than equal-step wins.

\emph{Limitations.} Fixed scatterers make long-horizon rates, and independent per-antenna gains absolute rates, optimistic; the findings hold with shared gains. The main results assume that the BS recomputes its precoder from the current channel with perfect knowledge. Pilot airtime (17--32\% lower rates), estimation error and 2--3-bit phases leave the slow solvers behind; at one bit a single-sweep search ties the best rounded Adam, and at pilot signal-to-noise ratios of 10~dB and below the fast methods cannot be told apart (Section~\ref{sec:res_robust}). RIS switching time is not modeled, and there is one geometry. Latencies come from single-threaded Python: compiled alone (62--63-fold faster), coordinate descent overtakes Python Adam, but an equal speed-up of all methods keeps tuned Adam ahead; the normalized step cost carries the results to other speeds and, approximately, carriers. Semidefinite relaxation~\cite{Wu2019TWC}, WMMSE~\cite{Choi2024TCOM} and manifold methods were not timed, and long-run Adam only at a fixed budget. In a first OFDM check the interior optimum held for centre-frequency optimization at 20~MHz, but at 50--100~MHz 48 steps were significantly worse in only 3 of 16 configurations.
